\documentclass[14pt,a4paper]{extarticle}

\usepackage[T2A]{fontenc}
\usepackage[utf8]{inputenc}
\usepackage[russian]{babel}
\usepackage[a4paper,margin=2cm]{geometry}
\usepackage{amsmath,amssymb}
\usepackage{graphicx}
\usepackage{booktabs}
\usepackage{float}
\usepackage{array}
\usepackage{caption}
\usepackage{hyperref}
\usepackage{setspace}

\onehalfspacing
\emergencystretch=3em
\graphicspath{{../outputs/lab4/plots/}}

\hypersetup{
    colorlinks=true,
    linkcolor=black,
    urlcolor=blue
}

\newcommand{\plotplaceholder}[2]{
    \begin{figure}[H]
        \centering
        \IfFileExists{../outputs/lab4/plots/#1}{
            \includegraphics[width=0.88\textwidth]{#1}
        }{
            \fbox{\parbox[c][5cm][c]{0.82\textwidth}{\centering Файл \texttt{\detokenize{#1}} создается запуском\\\texttt{python lab4/run.py}}}
        }
        \caption{#2}
    \end{figure}
}

\newcommand{\fitreporttable}[1]{
    \resizebox{\textwidth}{!}{#1}
}

\title{Лабораторная работа №4\\Сопряженные направления, методы Ньютона и квазиньютоновские методы}
\author{
    Новиков Алексей Георгиевич, группа М3233\\
    Душкин Александр Алексеевич, группа М3233
}
\date{2026}

\begin{document}

\begin{titlepage}
    \begin{center}
        {\large Отчет по лабораторной работе №4}\\[0.5cm]
        {\LARGE \textbf{Методы второго порядка и сопряженные направления}}\\[2cm]

        \begin{flushright}
            Выполнили:\\[0.2cm]
            Новиков Алексей Георгиевич, группа М3233\\
            Душкин Александр Алексеевич, группа М3233\\[0.5cm]
        \end{flushright}

        \vfill
        Санкт-Петербург\\[0.2cm]
        2026
    \end{center}
\end{titlepage}

\tableofcontents
\newpage

\section{Постановка задачи}

Цель работы --- сравнить методы, которые используют больше информации о геометрии функции, чем обычный градиентный спуск. В экспериментах рассматривались линейный метод сопряженных градиентов для квадратичных функций, нелинейные методы Флетчера--Ривса и Полака--Рибьера, методы Ньютона, Powell Dog Leg, DFP, BFGS, L-BFGS и библиотечный \texttt{Newton-CG} из SciPy.

Для каждого запуска фиксировались число итераций, количество вычислений функции, градиента и Гессиана, итоговое значение функции и факт сходимости по норме градиента. Основная точность:
\[
    \|\nabla f(x_k)\|\le 10^{-8}.
\]

\textbf{Промежуточный вывод.} В этой лабораторной сравниваются методы с разной ценой одной итерации. Сопряженные направления почти не используют вторые производные, методы Ньютона явно решают систему с Гессианом, а BFGS/DFP пытаются восстановить кривизну по истории шагов. Поэтому важно смотреть не только на итерации, но и на счетчики вызовов.

\section{Модели задач}

\subsection{Сгенерированные квадратичные функции}

Основная серия экспериментов проводилась на функциях
\[
    f(x)=\frac12 x^T A x - b^T x + c,
\]
где матрица $A$ симметричная положительно определенная. Размерность принимала значения $n\in\{2,10,50,100\}$, число обусловленности --- $k\in\{1,10,100,1000\}$, для каждой пары использовались два seed.

\textbf{Промежуточный вывод.} На квадратичных функциях хорошо видно, что именно дает учет кривизны. Если Гессиан известен и положительно определен, ньютоновское направление сразу указывает к минимуму. Если же метод строит приближение к обратному Гессиану по шагам, ему нужно несколько итераций, чтобы накопить информацию о форме уровней.

\subsection{Двумерные функции для траекторий}

Для визуализации использовались двумерная квадратичная функция, Розенброк, Химмельблау и Экли. Они нужны не столько для таблиц, сколько для проверки формы траекторий.

\textbf{Промежуточный вывод.} Таблица говорит, сколько шагов сделал метод, но не показывает, почему он двигался именно так. На графиках видно, когда метод идет к минимуму почти напрямую, когда срабатывают перезапуски или trust-region ограничения, а когда направление становится неудачным из-за нелинейной геометрии.

\section{Описание методов}

\subsection{Метод сопряженных градиентов для квадратичной задачи}

Для квадратичной функции с положительно определенной матрицей $A$ метод строит $A$-сопряженные направления:
\[
    p_{k+1}=-g_{k+1}+\beta_k p_k,\qquad
    \beta_k=\frac{g_{k+1}^Tg_{k+1}}{g_k^Tg_k}.
\]
В точной арифметике минимум находится не более чем за $n$ шагов. В вычислительных экспериментах это свойство может нарушаться из-за конечной точности, плохой обусловленности матрицы и строгого критерия по норме градиента, поэтому для таких запусков отдельно фиксируется факт сходимости.

\textbf{Промежуточный вывод.} Метод силен именно на квадратичных задачах: каждое новое направление не портит уже найденное уменьшение по предыдущим направлениям. Но это свойство опирается на фиксированную квадратичную форму и устойчивое сохранение сопряженности направлений; при плохой обусловленности и округлениях эта геометрия постепенно теряется. На нелинейных функциях такой гарантии уже нет.

\subsection{Нелинейные CG-методы}

Оба метода сохраняют идею сопряженных направлений, но пересчитывают коэффициент $\beta$ по текущим градиентам. В варианте Полака--Рибьера используется разность градиентов:
\[
    \beta_k^{PR}=\frac{g_{k+1}^T(g_{k+1}-g_k)}{g_k^Tg_k}.
\]
В реализации отрицательные значения обрезаются до нуля, что фактически делает перезапуск к антиградиенту.

\textbf{Промежуточный вывод.} Для нелинейных задач перезапуски важны: накопленное направление может перестать быть направлением спуска. Поэтому Polak--Ribiere+ часто устойчивее обычной формулы, но на сложных поверхностях все равно может тратить много итераций на поиск подходящего направления.

\subsection{Методы Ньютона и Powell Dog Leg}

Метод Ньютона решает систему
\[
    \nabla^2 f(x_k)p_k=-\nabla f(x_k)
\]
и затем делает шаг по найденному направлению. В версии с выбором направления Гессиан при необходимости модифицируется сдвигом, а если направление не является направлением спуска, используется запасной вариант. Powell Dog Leg работает в trust-region постановке: шаг выбирается внутри шара доверия между направлением Коши и ньютоновским шагом.

\textbf{Промежуточный вывод.} Ньютоновские методы быстры, когда Гессиан хорошо обусловлен и положительно определен. Но на нелинейных функциях Гессиан может быть неудачным локально, поэтому нужны защитные механизмы: сдвиг, выбор направления или trust-region радиус.

\subsection{DFP, BFGS и L-BFGS}

Квазиньютоновские методы не вычисляют Гессиан напрямую. Они обновляют приближение к обратному Гессиану по парам
\[
    s_k=x_{k+1}-x_k,\qquad y_k=g_{k+1}-g_k.
\]
BFGS использует устойчивое ранговое обновление, DFP --- альтернативное обновление обратной матрицы, а L-BFGS хранит только последние $m$ пар $(s,y)$ и применяет two-loop recursion.

\textbf{Промежуточный вывод.} Эти методы занимают середину между градиентными и ньютоновскими. Они дешевле по памяти и не требуют явного Гессиана, но качество направления зависит от того, насколько история шагов успела описать кривизну функции.

\section{Результаты}

\subsection{Средние метрики на квадратичных функциях}

\begin{table}[H]
\centering
\caption{Средние значения метрик на серии GeneratedQuadratic}
\fitreporttable{
\begin{tabular}{lrrrrr}
\toprule
Метод & Успешных запусков & Итерации & $N_f$ & $N_g$ & $N_H$\\
\midrule
NewtonCholesky & 32 & 1.00 & 3.00 & 3.00 & 1.00\\
NewtonDirectionChoice & 32 & 1.00 & 3.00 & 3.00 & 1.00\\
PowellDogLeg & 32 & 4.94 & 5.94 & 5.94 & 4.94\\
QuadraticConjugateGradient & 16 & 70.25 & 71.25 & 71.25 & 1.00\\
BFGS & 17 & 31.09 & 240.12 & 65.06 & 0.00\\
DFP & 16 & 27.97 & 200.28 & 58.31 & 0.00\\
LBFGS & 40 & 68.99 & 261.59 & 140.66 & 0.00\\
ScipyNewtonCG & 25 & 10.69 & 28.66 & 27.22 & 10.88\\
FletcherReeves & 8 & 239.91 & 2997.56 & 498.41 & 0.00\\
PolakRibiere & 8 & 148.78 & 1727.84 & 307.00 & 0.00\\
\bottomrule
\end{tabular}
}
\end{table}

В строке L-BFGS число успешных запусков считается по всем значениям памяти $m\in\{3,5,10\}$, поэтому оно может быть больше числа сгенерированных квадратичных задач. Для остальных методов доля успешных запусков показывает, сколько раз был выполнен строгий критерий $\|\nabla f(x_k)\|\le 10^{-8}$ до ограничения по итерациям.

\textbf{Промежуточный вывод.} Для чисто квадратичных задач методы Ньютона ожидаемо выигрывают по числу итераций: точный Гессиан полностью задает форму функции. Powell Dog Leg делает больше шагов, потому что trust-region радиус не всегда сразу разрешает полный ньютоновский шаг. Методы сопряженных градиентов и квазиньютоновские методы хуже переносят рост размерности и числа обусловленности: в конечной арифметике сопряженность направлений и положительность кривизны $s_k^Ty_k$ могут нарушаться, из-за чего появляются перезапуски, пропущенные обновления или остановка по лимиту итераций. Поэтому в полной таблице отдельно сохраняется сообщение остановки, а не только итоговое значение функции.

\subsection{Графики зависимостей}

\plotplaceholder{lab4_n_iter_vs_n_k_10.png}{Зависимость числа итераций от размерности при $k=10$}
\textit{Комментарий к графику.} У методов с явным использованием Гессиана рост размерности почти не увеличивает число итераций, но увеличивает стоимость линейной алгебры внутри шага. У методов без Гессиана число итераций сильнее реагирует на размерность, потому что информацию о кривизне приходится накапливать постепенно.

\plotplaceholder{lab4_n_iter_vs_k_n_10.png}{Зависимость числа итераций от числа обусловленности при $n=10$}
\textit{Комментарий к графику.} При росте $k$ уровни становятся вытянутыми. Методы, которые не учитывают кривизну напрямую, начинают делать больше корректирующих шагов, тогда как ньютоновские направления остаются ближе к оптимальному локальному шагу.

\plotplaceholder{lab4_comparison_n_iter.png}{Среднее число итераций: собственные методы и SciPy Newton-CG}
\textit{Комментарий к графику.} Сравнение с SciPy показывает не только скорость, но и цену универсальности. Библиотечный метод устойчиво использует Гессиан. Собственные методы проявляют свою специализацию: NewtonCholesky силен на SPD-квадратиках, L-BFGS экономит память, CG дешев по данным на итерацию.

\plotplaceholder{lab4_lbfgs_memory_n_iter.png}{Влияние размера памяти L-BFGS на число итераций}
\textit{Комментарий к графику.} Параметр $m$ задает, сколько последних пар $(s,y)$ участвует в восстановлении кривизны. Малое $m$ экономит память, но быстрее забывает геометрию задачи; большое $m$ дает больше информации для направления, но не всегда заметно улучшает результат.

\subsection{Траектории}

\plotplaceholder{Quadratic2DVisualization__seed_42_x0_-8.0_-6.0_x_star_1.0_-1.0__trajectories.png}{Траектории на двумерной квадратичной функции}
\textit{Комментарий к графику.} На квадратичной поверхности ньютоновское направление почти сразу попадает в минимум, а квазиньютоновским методам нужно несколько шагов, чтобы восстановить масштаб по осям.

\plotplaceholder{Lab4Rosenbrock__x0_-1.2_1.0__trajectories.png}{Траектории на функции Розенброка}
\textit{Комментарий к графику.} У Розенброка узкая изогнутая долина. Здесь полезна информация о кривизне, но слишком агрессивное направление требует линейного поиска или trust-region ограничения.

\plotplaceholder{Lab4Himmelblau__x0_2.0_2.0__trajectories.png}{Траектории на функции Химмельблау}
\textit{Комментарий к графику.} На Химмельблау важна область притяжения. После выбора локальной области методы различаются тем, насколько быстро и стабильно используют кривизну около минимума.

\plotplaceholder{Lab4Ackley__x0_1.0_1.0__trajectories.png}{Траектории на функции Экли}
\textit{Комментарий к графику.} У Экли много локальных колебаний. Методы второго порядка могут получить локально неудачный Гессиан, поэтому защитные механизмы и линейный поиск здесь важнее, чем на квадратике.

\section{Полная таблица экспериментов}

Полная сводка всех 552 запусков хранится в \texttt{../outputs/lab4/summary.csv}. Ниже приведена печатная версия той же таблицы, разбитая на масштабируемые блоки.

\input{../outputs/lab4/summary_report_scaled.tex}

\section{Общий вывод}

В этой работе хорошо видно разделение ролей между классами методов. Если задача квадратичная и Гессиан положительно определен, ньютоновский шаг практически сразу дает минимум. Если Гессиан недоступен или дорог, BFGS и L-BFGS становятся разумным компромиссом: они не требуют вторых производных, но постепенно восстанавливают информацию о кривизне.

Методы сопряженных градиентов дешевы по памяти и хорошо мотивированы на квадратичных задачах, однако на нелинейных функциях их качество сильнее зависит от линейного поиска и перезапусков. Trust-region подход Powell Dog Leg оказался более осторожным: он делает больше шагов, зато контролирует длину движения, когда локальная квадратичная модель может быть ненадежной.

\end{document}
